% GENERATED FILE. Edit canonical lessons, then run npm run generate:books.
% canonical-source-hash: fnv1a64:d4bb210382022f5f
% canonical-lessons: SA-C513-mathnati, SA-C513-manthah, SA-C513-ghrtam, SA-C513-galayati, SA-C513-snigdhah

\chapter{Milk and Butter}
\label{ch:sa-milk-and-butter}
\hlchaptermodality{513}

\begin{chapteropening}
\textbf{By the end of this chapter:} \emph{I can talk about churning, straining and ghee.}

Complete the last lesson of chapter 513: I can talk about churning, straining and ghee.
\end{chapteropening}

\section[\texorpdfstring{\sk{मथ्नाति} (mathnāti)}{मथ्नाति (mathnāti)}]{\sk{मथ्नाति} (mathnāti) — churns}

\label{lesson:SA-C513-mathnati}



\begin{quote}
Before the new one: say the Sanskrit for a winnowing basket, then the Sanskrit for combines.
\end{quote}

\subsection*{You'll want to know: \sk{मथ्नाति}}
\textbf{\sk{मथ्नाति}} — \emph{mathnāti} — \textquotedblleft{}churns\textquotedblright{}.

\begin{quote}
\sk{माता} \sk{दधि} \sk{मथ्नाति।}
\end{quote}

\begin{quote}
\emph{mātā dadhi mathnāti.}
\end{quote}

\begin{quote}
Mother churns the curd.
\end{quote}

Past (she …): \sk{सा} \sk{दधि} \sk{मथितवती}, \emph{sā dadhi mathitavatī}.

\begin{cousinweb}
\textbf{\sk{मथ्नाति}} (\emph{mathnāti}) comes from the root √math (also \emph{manth}), \textquotedblleft{}to stir, to churn\textquotedblright{}, whose present is \emph{mathnāti}; its past participle is \emph{mathita} (MW).
\end{cousinweb}

\begin{grammarlens}[title={What you've built}]
One more word for describing what you lived and felt.
\end{grammarlens}

\subsection*{Your turn}
\begin{itemize}
\raggedright
  \item \emph{Say it:} \emph{mathnāti}
  \item \emph{Say it:} \emph{mathnāti}, once more
  \item \emph{Say it:} \emph{mātā dadhi mathnāti.} — Mother churns the curd.
  \item \emph{Recall:} say the Sanskrit for a winnowing basket, then the Sanskrit for combines, then say \emph{mathnāti} again
\end{itemize}

\begin{tcolorbox}[breakable,colback=teal!4,colframe=teal!35!black,title={Before you move on}]
What does \sk{मथ्नाति} mean? (\textquotedblleft{}Churns\textquotedblright{}.) Say it once more.
\end{tcolorbox}
\section[\texorpdfstring{\sk{मन्थः} (manthaḥ)}{मन्थः (manthaḥ)}]{\sk{मन्थः} (manthaḥ) — a churning stick}

\label{lesson:SA-C513-manthah}



\begin{quote}
Before the new one: say the Sanskrit for combines, then the Sanskrit for churns.
\end{quote}

\subsection*{You'll want to know: \sk{मन्थः}}
\textbf{\sk{मन्थः}} — \emph{manthaḥ} — \textquotedblleft{}a churning stick\textquotedblright{}.

\begin{quote}
\sk{सा} \sk{मन्थेन} \sk{दधि} \sk{मथ्नाति।}
\end{quote}

\begin{quote}
\emph{sā manthena dadhi mathnāti.}
\end{quote}

\begin{quote}
She churns the curd with a churning stick.
\end{quote}

A masculine noun.

\begin{cousinweb}
\textbf{\sk{मन्थः}} (\emph{manthaḥ}) comes from the root √manth, \textquotedblleft{}to stir, to churn\textquotedblright{}; MW gives \emph{mantha} as churning and as a churning-stick (MW).
\end{cousinweb}

\begin{grammarlens}[title={What you've built}]
One more word for describing what you lived and felt.
\end{grammarlens}

\subsection*{Your turn}
\begin{itemize}
\raggedright
  \item \emph{Say it:} \emph{manthaḥ}
  \item \emph{Say it:} \emph{manthaḥ}, once more
  \item \emph{Say it:} \emph{sā manthena dadhi mathnāti.} — She churns the curd with a churning stick.
  \item \emph{Recall:} say the Sanskrit for combines, then the Sanskrit for churns, then say \emph{manthaḥ} again
\end{itemize}

\begin{tcolorbox}[breakable,colback=teal!4,colframe=teal!35!black,title={Before you move on}]
What does \sk{मन्थः} mean? (\textquotedblleft{}A churning stick\textquotedblright{}.) Say it once more.
\end{tcolorbox}
\section[\texorpdfstring{\sk{घृतम्} (ghṛtam)}{घृतम् (ghṛtam)}]{\sk{घृतम्} (ghṛtam) — ghee}

\label{lesson:SA-C513-ghrtam}



\begin{quote}
Before the new one: say the Sanskrit for churns, then the Sanskrit for a churning stick.
\end{quote}

\subsection*{You'll want to know: \sk{घृतम्}}
\textbf{\sk{घृतम्}} — \emph{ghṛtam} — \textquotedblleft{}ghee\textquotedblright{}.

\begin{quote}
\sk{नवनीतात्} \sk{घृतं} \sk{भवति।}
\end{quote}

\begin{quote}
\emph{navanītāt ghṛtaṃ bhavati.}
\end{quote}

\begin{quote}
Ghee is made from butter.
\end{quote}

A neuter noun.

\begin{cousinweb}
\textbf{\sk{घृतम्}} (\emph{ghṛtam}) comes from the root √ghṛ, \textquotedblleft{}to sprinkle\textquotedblright{}, as its past participle; MW gives \emph{ghṛta} as ghee, butter that has been boiled gently and left to cool (MW).
\end{cousinweb}

\begin{grammarlens}[title={What you've built}]
One more word for describing what you lived and felt.
\end{grammarlens}

\subsection*{Your turn}
\begin{itemize}
\raggedright
  \item \emph{Say it:} \emph{ghṛtam}
  \item \emph{Say it:} \emph{ghṛtam}, once more
  \item \emph{Say it:} \emph{navanītāt ghṛtaṃ bhavati.} — Ghee is made from butter.
  \item \emph{Recall:} say the Sanskrit for churns, then the Sanskrit for a churning stick, then say \emph{ghṛtam} again
\end{itemize}

\begin{tcolorbox}[breakable,colback=teal!4,colframe=teal!35!black,title={Before you move on}]
What does \sk{घृतम्} mean? (\textquotedblleft{}Ghee\textquotedblright{}.) Say it once more.
\end{tcolorbox}
\section[\texorpdfstring{\sk{गालयति} (gālayati)}{गालयति (gālayati)}]{\sk{गालयति} (gālayati) — strains}

\label{lesson:SA-C513-galayati}



\begin{quote}
Before the new one: say the Sanskrit for a churning stick, then the Sanskrit for ghee.
\end{quote}

\subsection*{You'll want to know: \sk{गालयति}}
\textbf{\sk{गालयति}} — \emph{gālayati} — \textquotedblleft{}strains\textquotedblright{}.

\begin{quote}
\sk{सः} \sk{क्षीरं} \sk{वस्त्रेण} \sk{गालयति।}
\end{quote}

\begin{quote}
\emph{saḥ kṣīraṃ vastreṇa gālayati.}
\end{quote}

\begin{quote}
He strains the milk through a cloth.
\end{quote}

Past (he …): \sk{सः} \sk{क्षीरं} \sk{गालितवान्}, \emph{saḥ kṣīraṃ gālitavān}.

\begin{cousinweb}
\textbf{\sk{गालयति}} (\emph{gālayati}) comes from the causative of the root √gal, \textquotedblleft{}to drip, to ooze\textquotedblright{}: MW gives \emph{gālayati} as to make drop, to filter, to strain (MW).
\end{cousinweb}

\begin{grammarlens}[title={What you've built}]
One more word for describing what you lived and felt.
\end{grammarlens}

\subsection*{Your turn}
\begin{itemize}
\raggedright
  \item \emph{Say it:} \emph{gālayati}
  \item \emph{Say it:} \emph{gālayati}, once more
  \item \emph{Say it:} \emph{saḥ kṣīraṃ vastreṇa gālayati.} — He strains the milk through a cloth.
  \item \emph{Recall:} say the Sanskrit for a churning stick, then the Sanskrit for ghee, then say \emph{gālayati} again
\end{itemize}

\begin{tcolorbox}[breakable,colback=teal!4,colframe=teal!35!black,title={Before you move on}]
What does \sk{गालयति} mean? (\textquotedblleft{}Strains\textquotedblright{}.) Say it once more.
\end{tcolorbox}
\section[\texorpdfstring{\sk{स्निग्धः} (snigdhaḥ)}{स्निग्धः (snigdhaḥ)}]{\sk{स्निग्धः} (snigdhaḥ) — oily; rich}

\label{lesson:SA-C513-snigdhah}



\begin{quote}
Before the new one: say the Sanskrit for ghee, then the Sanskrit for strains.
\end{quote}

\subsection*{You'll want to know: \sk{स्निग्धः}}
\textbf{\sk{स्निग्धः}} — \emph{snigdhaḥ} — \textquotedblleft{}oily; rich\textquotedblright{}.

\begin{quote}
\sk{घृतेन} \sk{ओदनः} \sk{स्निग्धः} \sk{भवति।}
\end{quote}

\begin{quote}
\emph{ghṛtena odanaḥ snigdhaḥ bhavati.}
\end{quote}

\begin{quote}
Ghee makes the rice rich.
\end{quote}

\begin{cousinweb}
\textbf{\sk{स्निग्धः}} (\emph{snigdhaḥ}) comes from the root √snih, \textquotedblleft{}to be sticky, to be fond of\textquotedblright{}, as its past participle; MW gives \emph{snigdha} as sticky, unctuous, smooth (MW).
\end{cousinweb}

\begin{grammarlens}[title={What you've built}]
One more word for describing what you lived and felt.
\end{grammarlens}

\subsection*{Your turn}
\begin{itemize}
\raggedright
  \item \emph{Say it:} \emph{snigdhaḥ}
  \item \emph{Say it:} \emph{snigdhaḥ}, once more
  \item \emph{Say it:} \emph{ghṛtena odanaḥ snigdhaḥ bhavati.} — Ghee makes the rice rich.
  \item \emph{Recall:} say the Sanskrit for ghee, then the Sanskrit for strains, then say \emph{snigdhaḥ} again
\end{itemize}

\begin{tcolorbox}[breakable,colback=teal!4,colframe=teal!35!black,title={Before you move on}]
What does \sk{स्निग्धः} mean? (\textquotedblleft{}Oily; rich\textquotedblright{}.) Say it once more.
\end{tcolorbox}
